%% Chapter 8 — Conclusions
\section{Conclusions and Outlook}
\label{sec:conclusions}

\subsection{Summary}

This thesis has documented the setup, characterisation, and initial evaluation of a
\QICK{}-based readout architecture on a ZCU216 RFSoC, together with the design of a
machine-learning-assisted discrimination pipeline on existing experimental data.
The main contributions are:

\begin{enumerate}
  \item \textbf{System integration}: The ZCU216 board, XM655 analogue front-end,
        and \QICK{} firmware were set up for use in data collection, including SSH
        access, PYNQ configuration, and a Jupyter-notebook workflow for running
        experiments on the board.

  \item \textbf{Phase-coherent readout calibration}: Following the \QICK{} examples,
        the phase-calibration procedure was set up and validated on this system: it
        measures the random DAC--ADC DDS phase offset at each power-on and fits the
        linear model $\phi(f) = \phi_0 - 360^\circ \tau f$, with the fitted
        \texttt{res\_phase} folded into the standard experiment configuration.

  \item \textbf{Readout data-path}: The decimated and accumulated readout modes
        were characterised.  Accumulated readout gives one integrated $(I, Q)$ pair
        per shot for SSRO; repeat-averaging improves spectroscopy SNR but destroys
        the single-shot information.  A resonator spectroscopy workflow was
        implemented and validated in loopback; its extension to multiplexed
        readout was assessed but not implemented.

  \item \textbf{Real-time feedback primitives}: The \QICK{} \texttt{condj}
        conditional-branch demo was reproduced in loopback, and an active-reset
        sequence (dispersive readout, thresholding, calibrated $\pi$ pulse) was
        outlined from these primitives, with an estimated digital decision-to-fire
        latency below \SI{100}{\nano\second}.

  \item \textbf{ML-assisted discrimination}: On integrated IQ data, LDA reaches an
        assignment fidelity of 90.6\% ($\ket{0}$: 94.0\%, $\ket{1}$: 87.3\%), and a
        grid search over 20 MLP architectures finds no nonlinear classifier that
        improves on it, as expected for near-Gaussian, linearly separable blobs.
        The $\ket{1}$ deficit is the signature of $T_1$ relaxation mid-shot.  A
        953-parameter 1D CNN on synthetic raw traces reaches 99.75\% overall
        fidelity and 98.4\% accuracy on the 122 relaxation shots, against 44.3\%
        for LDA; because the synthetic set over-represents relaxation, the robust
        conclusion is the qualitative advantage of raw-trace input rather than the
        size of the gap.  FPGA deployment remains future work;
        Section~\ref{subsec:deployment_path} outlines the path toward it.
\end{enumerate}

\subsection{Outlook}

Several extensions follow directly from the results above.  The first experimental
priority is to collect a real \texttt{SSRO\_raw.h5} dataset from the ZCU216 with
\texttt{acquire\_decimated}, retrain the CNN, and check that the $T_1$-relaxation
recovery generalises to real data; the synthetic $T_1$ should first be recalibrated
to the measured $\ket{1}$ error rate of 12.7\%, an upper bound on the relaxation
fraction (part of it is Gaussian-overlap error, cf.\ the 6.0\% $\ket{0}$ error
rate), implying $T_1 \gtrsim 1500$ samples rather than the 600 used
here.  With real traces in hand, the
CNN can be exported via \texttt{hls4ml} and synthesised as a fixed-point IP block in
the QICK bitstream (Section~\ref{subsec:deployment_path}), closing the active-reset
loop at full speed without a software round-trip.

Beyond a single qubit, extending to multiplexed readout via the \QICK{} PFB firmware
would allow simultaneous readout of all qubits in Qilimanjaro's multi-qubit chips, a
prerequisite for quantum error-correction cycles.  The \QICK{}-based readout layer
should also be integrated with Qilimanjaro's higher-level programming interfaces
(circuit compilation, calibration databases, experiment management) for a
production-ready stack.  Finally, the fluxonium dispersive shift varies strongly
across the external-flux landscape, and an ML-assisted optimisation of the readout
circuit parameters to maximise the cavity pull over that landscape is a natural
complementary direction.
